\chapter[The cross-level identity (2) <-> (3): Pfaffian-Dirac, orientation character, primitive recognition]{The cross-level identity \protect\textcircled{\scriptsize 2}\,\(\leftrightarrow\)\,\protect\textcircled{\scriptsize 3}: Pfaffian--Dirac, orientation character, primitive recognition}
\label{ch:pfaffian-dirac}

The Igusa cusp form \(\Delta_5\) of weight five on
\(\Sp(2,\Z)\backslash\mathcal H_2\) carries three names, each rooted in a
distinct categorical-dimensional level:

\begin{itemize}
\item \(\mathcal D_X=\Delta_5\) — the Borcherds--Gritsenko--Nikulin section of the
  weight-\(5\) automorphic line over the genus-two Siegel space (View
  \textcircled{\scriptsize 1});
\item \(\operatorname{den}(\mathfrak g_{\Delta_5})=64^{-1}\Delta_5(2Z)\) —
  the denominator of the generalized Borcherds--Kac--Moody (BKM)
  superalgebra \(\mathfrak g_{\Delta_5}\) on
  \(\La^{2,1}_{II}\) (View \textcircled{\scriptsize 2});
\item \(\operatorname{Pf}_{\mathrm{prot}}(\mathfrak D_X^{\mathrm{DI}})\) —
  the protected Pfaffian of the Dirac--Igusa pro-object on
  \(X=K3\times E\) (View \textcircled{\scriptsize 3}).
\end{itemize}

The cross-level identity
\textcircled{\scriptsize 1}\(\leftrightarrow\)\textcircled{\scriptsize 2}
is Borcherds 1995 and Gritsenko--Nikulin 1998.  The edge
\textcircled{\scriptsize 2}\(\leftrightarrow\)\textcircled{\scriptsize 3}
has two pieces.  The four obstruction classes and the finite geometric
Pfaffian datum \((P_{\mathrm{fin}})\) give the Pfaffian; the
orientation classes give the orientation character on \(K3\times E\).
The Lie-superalgebra recognition has the additional finite
compact-source datum \((S1)\)--\((S10)\) and the finite Koszul
comparison datum of
Definition~\ref{def:finite-koszul-comparison-datum}.  The three
statements have separate hypotheses:

\begin{enumerate}
\item the \emph{Pfaffian--Dirac theorem}: under (D0), (O1), (O1)\(^+\),
  (O2), and the finite geometric Pfaffian datum
  \((P_{\mathrm{fin}})\) on a cofinal HN system,
  \[
    \operatorname{Pf}_{\mathrm{prot}}(\mathfrak D_X^{\mathrm{DI}})
    \;=\;\Delta_5
  \]
  as sections of \(\mathcal L^5\otimes\nu_{\Delta_5}\);
\item the \emph{orientation character identity}: under (O1)+(O1)\(^+\)+(O2),
  \[
    \epsilon_o\;=\;\nu_{\Delta_5}
    \quad\text{on}\quad W^{(2)}(\La^{2,1}_{II});
  \]
\item the \emph{primitive recognition theorem}: under the finite
  compact-source recognition datum (Cartan, simple representatives,
  Chevalley, real Serre, imaginary orthogonality, generation,
  completed PBW, Hopf radical equality, no-extra-relations, full parity
  dimensions),
  \[
    H^0\Prim_{\mathrm{prot}}\!\bigl(
    \overline\Pi_{X,*}^{\Theta}\mathcal F_X\bigr)
    \big/\overline{\operatorname{Rad}\langle\,\cdot\,,\,\cdot\,\rangle_X}
    \;\cong\;\mathfrak g_{\Delta_5}
  \]
  after passage to the inverse limit over the cofinal root windows, with
  the chiral cobar
  \(\Omega_E^{\mathrm{ch}}C_X\) identified as the formal current envelope
  \(\mathsf{Den}_{\Delta_5,E}\) only after the chiral Koszul source
  datum is also fixed.
\end{enumerate}

The theorem-dependency ledger for these three assertions is
\[
\mathrm{beilinson\_levels},\quad
\mathrm{objects},\quad
\mathrm{morphisms},\quad
\mathrm{theorem\_dependencies},\quad
\mathrm{dependency\_edges},\quad
\mathrm{forbidden\_edges},\quad
\mathrm{status\_separation},\quad
\mathrm{level\_equalities},\quad
\mathrm{equality\_sites},\quad
\mathrm{definition\_separation},\quad
\mathrm{obstruction\_independence},\quad
\mathrm{transitions},\quad
\mathrm{scalar\_firewall}.
\]
It is the finite graph recording which retained packets enter each
proof and which edges are forbidden.  In particular, the Pfaffian and
orientation proofs do not use the mirror conjecture, the gravity-line
residual, the scalar trace, or the OP scalar branch; primitive
recognition does not use the scalar trace, the denominator product, a
target counit, or target labels as source representatives.  Claim
statuses do not collapse: imported BKM and automorphic identities,
retained-data Pfaffian identities, conditional recognition, mirror
conjectures, scalar trace identities, and open gravity-line residuals
occupy distinct rows.
The definition-separation rows assert non-circularity at the level of
finite data: the Pfaffian line is not defined by the automorphic
section or by its desired value, the orientation character is not
defined by scalar signs, the chiral Koszul source is not the target
counit, the primitive source is not defined by target multiplicities,
and the Dirac--Igusa pro-object is not defined by a scalar shadow.
The obstruction-independence rows assert that \((D0)\), \((O1)\),
\((O1)^+\), \((O2_{\mathrm{atlas}})\), \(P_{\mathrm{fin}}\), the
Koszul source comparison, and primitive recognition are separate finite
data; typing dependencies are not logical implications between
vanishing statements.
The certificate
\(\texttt{certificates/dependencies/k3e\_theorem\_dependencies}\)
is schema-complete only: its verifier returns
\(\texttt{SCHEMA\_COMPLETE\_SCHEMA\_ONLY}\) and records
\(\texttt{dependency\_certification=false}\) and
\(\texttt{mathematical\_certification=false}\).  It verifies that the
finite ledger is typed, populated, separated by status, and guarded by
forbidden-edge, definition-separation, and obstruction-independence
rows; it does not prove the retained Pfaffian,
orientation, Koszul, recognition, mirror, scalar-trace, or gravity-line
claims.

\begin{proposition}[Hypothesis-deletion residues]
\label{prop:ch6-hypothesis-deletion-residues}
The theorem lane has the following exact residues when one named datum
is omitted.
\begin{enumerate}
\item Without the finite compact-source recognition datum
\((S1)\)--\((S10)\), the Pfaffian--Dirac section identity and the
orientation character identity remain under their own hypotheses, but
there is no theorem identifying
\[
P_X^{\Pi,\mathrm{red}}\quad\text{with}\quad\mathfrak g_{\Delta_5}.
\]
The surviving primitive object is only the retained Hall primitive
object with its supplied product, coproduct, pairing, radical, and PBW
rows; denominator exponents or signed multiplicities do not determine
its Lie bracket.

\item Without the chiral Koszul source datum, the finite Koszul
comparison datum, and strict Koszul Mittag--Leffler exactness, the
source coalgebra \(C_X\) is not identified with the target current
envelope.  The theorem no longer supplies a quasi-isomorphism
\[
\Omega_E^{\mathrm{ch}} C_X
\longrightarrow
\mathsf{Den}_{\Delta_5,E}
\]
with acyclic cone.  The exact residue is the compact
Hall primitive source, if supplied, before the source-to-target chiral
comparison.

\item Without the finite geometric Pfaffian datum
\((P_{\mathrm{fin}})\), the obstruction data may still define
orientation and local wall signs, but there is no global object
\[
\operatorname{pf}_{X,R}\in H^0(\mathfrak Q_R,\mathscr L_{\mathrm{Pf},R})
\]
whose inverse limit can be compared with \(\mathcal D_X=\Delta_5\).
The residue is a reduced orientation/wall datum, not a Pfaffian section
identity.

\item Without \((O2_{\mathrm{atlas}})\), the retained D0-HN datum,
quotient orientation, and Weyl-lift torsor do not compute the local
rank-one Pfaffian signs.  The residue is a Weyl-transported
orientation line on the retained quotient, but no theorem gives
\(\epsilon_o(s_{\delta_i})=-1\), no divisor order along the type-II
walls is computed, and the Pfaffian section is not identified with
\(\Delta_5\).

\item Without \((O1)^+\), the quotient orientation may exist, but it is
not transported by a coherent action of
\(W^{(2)}(\Lambda_{II}^{2,1})\).  The residue is an oriented reduced
quotient on the retained finite stages; there is no global character
\[
\epsilon_o:W^{(2)}(\Lambda_{II}^{2,1})\longrightarrow\{\pm1\}.
\]

\item Without \((O1)\), the quotient square root and reduced
orientation line on the \(K3\times E\)-quotient self-Ext frame bundle
are absent.  The residue is the finite d-critical/cosection and
D0-specialization data before orientation; Hall products and Pfaffian
sections are unsigned and do not define the protected oriented
Pfaffian.

\item Without \((D0)\), the cosection-reduced orientation theory has no
D0-degeneration endpoint on the cofinal HN pro-system.  The residue is
the formal Mukai--Gram lattice, the imported automorphic/BKM/theta
triangle, and any finite stage objects supplied independently; it is
not a compact \(K3\times E\) Pfaffian--Dirac theorem.
\end{enumerate}
Each residue is a positive statement about the surviving object.  None
of the rows permits a scalar trace, an automorphic target, or a
denominator product to replace the deleted datum.
\end{proposition}

\begin{proof}
The conclusion is a direct reading of the source and target of each
comparison.  Primitive recognition needs the source representatives,
relations, radical, pairing, parity, and PBW rows; without them there
is no Lie-superalgebra isomorphism.  The Koszul comparison is the only
map from the source coalgebra to the target current envelope.  The
finite geometric Pfaffian datum supplies the Pfaffian line and section,
so deleting it removes the domain of the section identity.  The
\((O2_{\mathrm{atlas}})\) datum supplies the local rank-one signs at
the type-II walls; deleting it removes the generator values of the
orientation character.  The \((O1)^+\) datum supplies coherent Weyl
transport; deleting it removes the homomorphism from the Weyl group.
The \((O1)\) datum supplies the quotient orientation line; deleting it
removes the oriented Pfaffian input.  The \((D0)\) datum supplies the
D0-degeneration endpoint and transition-compatible reduced orientation
theory; deleting it removes the compact pro-limit in which the
Pfaffian--Dirac comparison lives.  These are object-level failures in
the Beilinson chain, not failures of notation.
\end{proof}

\begin{proposition}[Implication and non-redundancy of obstruction rows]
\label{prop:ch6-obstruction-independence}
Let
\[
\mathsf{Obs}^{\mathrm{DI}}_{\mathrm{fin}}
=
\mathsf{D0}\times\mathsf{O1}\times\mathsf{O1}^+
\times\mathsf{O2}\times\mathsf{Pfin}\times
\mathsf{Kos}\times\mathsf{Rec}
\]
be the finite theorem-dependency product whose seven factors are the
row categories for \(D0_{\mathrm{HN}}\), \(O1_{\mathrm{quot}}\),
\(\mathcal T_W\), \(O2_{\mathrm{atlas}}\), \(P_{\mathrm{fin}}\), the
chiral Koszul comparison datum, and the compact-source primitive
recognition datum.  Each vanishing predicate is a predicate on its own
factor.  The only implication used by the Pfaffian--Dirac lane is the
conjunction of the relevant vanishing predicates into the corresponding
theorem:
\[
D0\wedge O1\wedge O1^+\wedge O2\wedge P_{\mathrm{fin}}
\Longrightarrow
\text{Pfaffian section and orientation character},
\]
and, separately,
\[
\text{compact source}\wedge\text{Koszul comparison}\wedge
\text{primitive recognition}
\Longrightarrow
P_X^{\Pi,\mathrm{red}}\cong\mathfrak g_{\Delta_5}.
\]
There is no theorem-level implication from any one named vanishing
predicate to any other.  The type dependencies recorded in
\(\texttt{obstruction\_independence.csv}\) say only that a later datum
is not meaningful until its earlier input object exists; they do not
prove the earlier datum.

For each of the seven rows, deleting that row and leaving the other row
slots formally zero gives a finite theorem-dependency witness to
non-redundancy.  The residual theorem in the deleted row is exactly the
corresponding item of
Proposition~\ref{prop:ch6-hypothesis-deletion-residues}.  These are
counterexamples to illegitimate implications inside the finite
dependency calculus.  They are not asserted to be geometric
\(K3\times E\) families with all other analytic hypotheses realized.
\end{proposition}

\begin{proof}
In the product category
\(\mathsf{Obs}^{\mathrm{DI}}_{\mathrm{fin}}\), a vanishing predicate is
pulled back from one projection.  A predicate pulled back from one
projection cannot force a predicate pulled back from another projection
unless such an implication is added as a row of the dependency ledger.
The table
\(\texttt{obstruction\_independence.csv}\) requires
\(\texttt{logical\_implications=none}\) for the seven named data and
records the type dependencies separately.  Thus the row \(O1^+\) being
typed over \(O1\) means that Weyl transport needs an orientation line;
it is not a proof that the orientation line exists.  The same applies
to \(O2_{\mathrm{atlas}}\), \(P_{\mathrm{fin}}\), the Koszul
comparison, and primitive recognition.

Deleting one row removes the source, target, or comparison named in
Proposition~\ref{prop:ch6-hypothesis-deletion-residues}, while leaving
the remaining row slots available as formal data.  This supplies the
finite-row countermodel to every attempted proof that the deleted row
is redundant.  Since each deleted row removes a different object or
comparison map, no two rows have the same residue; the obstruction
classes are non-redundant in the theorem-dependency calculus.
\end{proof}

The eight finite constructions of \S5.7 define, at every finite stage \(R\)
in the Harder--Narasimhan inverse system, the data \(\mathcal A_{X,R}^{E_3}\),
\(F_{X,R}^{\mathrm{hyb}}\), \(\Gamma_{X,R}\), \(\Pi_X\), \(\Phi_R\),
\(o_R\), \(H_R\), \(P_R^\Pi\), \(C_{X,R}\), \(\Theta_{\mathrm{Kos},R}\),
\(\mathscr L_{\mathrm{Pf},R}\), \(\operatorname{pf}_{X,R}\),
\(\varepsilon_{o,R}\), and \(\mathrm{Rec}_R\) of the pro-object
\(\mathfrak D_X^{\mathrm{DI}}\); the four obstruction classes of \S5.8
are exactly (D0), (O1), (O1)\(^+\), (O2).  These four classes act on
the Pfaffian line and the orientation character; the assertion that
\(\mathrm{Rec}_R\) is eventually an isomorphism is the separate
compact-source recognition datum.  The relevant geometric data are the
open factorization chart of the vacuum \(b\), the protected Pfaffian
functor, the Hall--Drinfeld correspondence, the chosen
Sp\(^{\mathrm{ch}}\) specialization, the pro-object ambient, the
D0-degeneration endpoint, the Bussi--Brav--Dupont--Joyce--Szendr\H{o}i
orientation, the non-degeneracy of the Pfaffian line, and the type-II
wall normal form.

The Pfaffian and
orientation theorems are retained-data theorems on \(K3\times E\), as in
Appendix~\ref{app:obstruction-discharges}; the primitive recognition
theorem is conditional on the compact-source datum
\((S1)\)--\((S10)\).  The BKM denominator identity (View
\textcircled{\scriptsize 2}) is
\emph{proved}, imported from
\cite{Borcherds95,GN,GNII,BorcherdsGKM88,KAC:1}; the automorphic
identification (View \textcircled{\scriptsize 1}) is \emph{proved},
imported from \cite{Igusa1962,Igusa1964II}; the operator-level
Pfaffian datum \(\mathfrak D_X^{\mathrm{DI}}\) is \emph{proved on
\(K3\times E\) relative to the retained D0-HN datum, the retained
quotient-orientation datum, the chosen Weyl lifts,
\((O2_{\mathrm{atlas}})\), and the retained finite geometric Pfaffian datum
\((P_{\mathrm{fin}})\)}
of
Appendix~\ref{app:obstruction-discharges}; the chiral cobar
identification is \emph{conditional} on the Koszul source datum of
Definition~\ref{def:chiral-koszul-source-datum}.  The
\(\kappa\)-tuple of the Calabi--Yau quantum-groups companion,
\(
(\kappa_{\mathrm{cat}},\kappa_{\mathrm{ch}}^{\mathrm{Hodge}},
\kappa_{\mathrm{ch}}^{\mathrm{Heis}},\kappa_{\mathrm{BKM}},
\kappa_{\mathrm{fiber}})=(0,0,3,5,24),
\)
fixes the level of \(\Delta_5\): the cusp-form weight is
\(\kappa_{\mathrm{BKM}}=5\) and the elliptic genus weight is
\(\chi_{\mathrm{top}}(K3)=\kappa_{\mathrm{fiber}}=24\); the additive
formula \(\kappa_{\mathrm{BKM}}=\kappa_{\mathrm{ch}}+\chi(\mathcal O_{\mathrm{fiber}})\)
has no invariant meaning on the \(\kappa\)-tuple.


\section{Data and Obstructions}
\label{sec:ch6-data}

\subsection{The pro-object \(\mathfrak D_X^{\mathrm{DI}}\) and its eight finite-stage components}
\label{subsec:ch6-pro-object}

Three names of \(\Delta_5\): the section \(\mathcal D_X=\Delta_5\) (View
\textcircled{\scriptsize 1}); the BKM denominator
\(\operatorname{den}(\mathfrak g_{\Delta_5})=64^{-1}\Delta_5(2Z)\) (View
\textcircled{\scriptsize 2}); the protected Pfaffian
\(\operatorname{Pf}_{\mathrm{prot}}(\mathfrak D_X^{\mathrm{DI}})\) (View
\textcircled{\scriptsize 3}, relative to retained data).  The handle
``\(\Delta_5\)'' is reserved for the working symbol; theorem-level
sentences specify the named view.

The Dirac--Igusa pro-object is
\[
\mathfrak D_X^{\mathrm{DI}}
\;=\;\lim_R\bigl(
\mathcal A_{X,R}^{E_3},\,
F_{X,R}^{\mathrm{hyb}},\,
\Gamma_{X,R},\,
\Pi_X,\,
\Phi_R,\,
o_R,\,
H_R,\,
P_R^\Pi,\,
C_{X,R},\,
\Theta_{\mathrm{Kos},R},\,
\mathscr L_{\mathrm{Pf},R},\,
\operatorname{pf}_{X,R},\,
\varepsilon_{o,R},\,
\mathrm{Rec}_R\bigr).
\]
The components are the values of the eight finite constructions of \S5.7,
with the ambient category the cofinal HN inverse system of
Definition~\ref{def:normal-ordered-charge-lattice-appendixC}; \(R\)
ranges over a chosen cofinal subsystem
\(\mathcal R_{\mathrm{HN}}\subset\Lambda_{\mathrm{ample}}\).  The
transition maps are the finite-stage morphisms of
Definition~\ref{def:finite-stage-dirac-igusa-morphism}; isomorphism of
two such pro-objects means pro-isomorphism in the sense of
Definition~\ref{def:dirac-igusa-pro-isomorphism}.  At finite stage,
these transition and pro-isomorphism claims are the
\(\mathrm{cofinal\_subsystems}\), \(\mathrm{stage\_objects}\),
\(\mathrm{stage\_morphisms}\), \(\mathrm{component\_maps}\),
\(\mathrm{composition\_laws}\), \(\mathrm{ml\_exactness}\),
\(\mathrm{pro\_isomorphisms}\), and \(\mathrm{scalar\_firewall}\)
tables of \S\ref{subsec:eight-finite-constructions}.  The label
\(\gamma\) is the pro-object ambient.
The universal property of this pro-object is the Hom formula of
Proposition~\ref{prop:dirac-igusa-pro-object-universal-property}; in
particular, changing the finite HN presentation by a cofinal subsystem
does not change \(\mathfrak D_X^{\mathrm{DI}}\) in the pro-category.

\begin{remark}[The three \(\Delta_5\)-names]
\label{rem:ch6-three-delta-names}
The equality
\(\operatorname{Pf}_{\mathrm{prot}}(\mathfrak D_X^{\mathrm{DI}})=\mathcal D_X\) is a
cross-level identity, not a tautology.  The right-hand side is an
automorphic section; the left-hand side is the inverse limit of finite
Pfaffian sections \(\operatorname{pf}_{X,R}\) of the line bundles
\(\mathscr L_{\mathrm{Pf},R}\) on the reduced quotient stacks
\(\mathfrak Q_R\); the two are related by the fixed comparison
\(\iota_{\mathrm{aut}}:\mathscr L_{\mathrm{Pf},X}\xrightarrow{\sim}
\mathcal L^5\otimes\nu_{\Delta_5}\).  The denominator name
\(\operatorname{den}(\mathfrak g_{\Delta_5})=64^{-1}\Delta_5(2Z)\) is
the same global section \(\mathcal D_X\) written in the type-II chamber
polarization that exposes the Weyl--Kac--Borcherds product expansion;
the prefactor \(64^{-1}\) and the doubling \(2Z\) of the modular
variable enter from the theta normalization of the Borcherds product
on the index-$4$ sublattice $\Lambda^{2,1}_{II}\subset\Lambda^{2,1}$.
The source side \(P_X^{\Pi,\mathrm{red}}\) identifies with
\(\mathfrak g_{\Delta_5}\) only after the recognition criterion of
\S\ref{sec:ch6-primitive-recognition} is established.
\end{remark}

\subsection{The four obstruction classes: (D0), (O1), (O1)\(^+\), (O2)}
\label{subsec:ch6-obstructions}

The four classes match the eight finite constructions one-to-one only on the
orientation lane.  They are the obstruction data of
Theorem~\ref{thm:pfaffian-dirac-finite-stage}.  On \(K3\times E\), Appendix~\ref{app:obstruction-discharges}
reduces \((D0)\) to the retained D0-HN datum by
Theorem~\ref{thm:G-D0}; reduces \((O1)\) to the
retained quotient-orientation datum by Theorem~\ref{thm:G-O1}; reduces
\((O1)^+\) to the chosen Weyl lifts by Theorem~\ref{thm:G-O1-plus};
and reduces \((O2)\) to the retained type-II wall-atlas datum
\((O2_{\mathrm{atlas}})\) by Theorem~\ref{thm:G-O2-orbit-transport}.  The
\((R1)\)--\((R2)\) scalar separation for the type-II middle-wall
constant \(4096=2^{12}=64^2\) inside \((O2)\) is
Theorem~\ref{thm:G-R1-R2}.
The finite D0 packet
\(\texttt{certificates/d0/k3e\_d0\_hn}\) is presently verified only
through its obstruction ledger:
\(\texttt{D0\_OBSTRUCTION\_LEDGER\_VERIFIED}\), with
\(\texttt{d0\_certification=false}\).  Hence the D0-dependent
assertions in this chapter are relative theorems until the flat
degeneration, retained-substack, HN-transition, derived-enhancement,
semiregularity-cosection, vanishing-cycle, orientation, Pfaffian,
protected-integration, and Mittag--Leffler rows are supplied.
The finite geometric Pfaffian packet
\(\texttt{certificates/pfaffian/k3e\_finite\_pfaffian}\) is likewise
verified only through its obstruction ledger:
\(\texttt{PFAFFIAN\_OBSTRUCTION\_LEDGER\_VERIFIED}\), with
\(\texttt{pfaffian\_certification=false}\).  Hence the protected
Pfaffian identity below is a relative section identity until the skew
perfect complexes, Pfaffian lines, sections, wall charts,
automorphic-comparison rows, and Pfaffian-line transition rows are
constructed.
The finite O2 wall-atlas packet
\(\texttt{certificates/wall\_atlas/k3e\_o2\_atlas}\) is also verified
only through its obstruction ledger:
\(\texttt{O2\_OBSTRUCTION\_LEDGER\_VERIFIED}\), with
\(\texttt{o2\_certification=false}\).  Hence every assertion using
\((O2_{\mathrm{atlas}})\) is relative until the simple-wall coverage,
wall-object, wall-chart, overlap, Weyl-transport, half-Hilbert-orbit,
compatibility, type-I-exclusion, D0-limit, and scalar-separation rows
are supplied.
The finite orientation packet
\(\texttt{certificates/orientation/k3e\_reduced\_orientation}\)
is presently verified only through its obstruction ledger:
\(\texttt{ORIENTATION\_OBSTRUCTION\_LEDGER\_VERIFIED}\), with
\(\texttt{orientation\_certification=false}\).  Thus the
orientation-dependent assertions in this chapter are relative
theorems until the square-root, quotient Cech--Borel,
finite-stabilizer, Weyl-lift, Coxeter, and transition rows are supplied.

\paragraph{(D0) — D0-degeneration endpoint of the cosection-reduced
d-critical orientation theory.}  At every
\(R\), the retained \(D0_{\mathrm{HN}}\) datum of
Definition~\ref{def:G-D0-HN-system} is fixed: a flat
D0-degeneration family, retained finite-type object/extension/flag
stacks, proper or closed HN transition morphisms, quasi-smooth
d-critical derived enhancements, surjective additive K3 semiregularity
cosections, reduced vanishing-cycle specialization, Joyce--Upmeier
orientation specialization, Pfaffian-line and protected-integration
specialization, and strict Mittag--Leffler exactness.  Compatibly in
\(R\), these clauses give the HN-completed state object and the finite
first-order block \(\mathfrak D_X\).
The finite first-order block also carries the exponent-matching clause
\[
\sdim P^\Pi_{X,(n,l,m)}=f(nm,l)
\qquad ((n,l,m)\in\Gamma_{\mathrm{eff}}),
\]
where \(f(n,l)\) is the Fourier coefficient of the weak Jacobi form
\(\phi_{0,1}\) of weight zero and index one.  This clause is not a
consequence of the reduced orientation alone; it is part of the
retained finite-stage Dirac block.  The cosection contraction
of the Pantev--Toen--Vaquie--Vezzosi (PTVV)
\((-1)\)-shifted symplectic form by the K3 semiregularity cosection
\(\sigma_{K3}\) is the Maulik--Toda / Oberdieck reduced obstruction theory
\cite{MaulikTodaGV,OberdieckReducedPT}; this absorbs the parity flip
\(\mathrm{vd}^{\mathrm{red}}=\mathrm{vd}^{\mathrm{std}}+1\) into the
reduced orientation line as the \(\Z/2\)-grading shift
\(o^{\mathrm{red}}_{\mathcal C}=o^{\mathrm{std}}_{\mathcal C}\otimes\Pi\).

\paragraph{(O1) — strong multiplicative reduced orientation datum on the
reduced \(K3\times E\)-quotient self-Ext frame bundle.}  A Joyce--Upmeier-type orientation line
\(o_{\mathcal C}\in\operatorname{Pic}(\mathfrak M^{or}_X)\) with
\(o_{\mathcal C}^{\otimes 2}\simeq
\det\operatorname{RHom}_{\mathrm{red}}(\mathcal C,\mathcal C)\),
extension-multiplicative under
\(o_{\mathcal C\oplus\mathcal C'}\simeq
o_{\mathcal C}\otimes o_{\mathcal C'}\), and equipped with the full
quotient-orientation datum
\[
\theta_{\mathcal C}=
\bigl(\theta^{\mathrm{red}}_{\mathcal C},\theta^{E,\mathrm{free}}_{\mathcal C},
\{(\theta^{H}_{\mathcal C},\lambda^{H}_{\mathcal C}=0)\}_{H\in\mathcal H(\mathcal C)}\bigr)
\]
on every retained charge stratum.  The vanishings
\(\alpha^{\mathrm{red}}_{\mathcal C}=0\),
\(\alpha^{E,\mathrm{free}}_{\mathcal C}=0\),
\(\widetilde\beta^{H}_{\mathcal C}=0\),
\(\lambda^{H}_{\mathcal C}=0\) are the orientation-line, the Borel-equivariant
free-action, the residual finite-stabilizer Borel, and the residual
finite-stabilizer linearization conditions; for \(N=2\), the Klein-four
calculation
\(\beta^{E,E[2]}_{\mathcal C}=b_{20}x_1^2+b_{11}x_1x_2+b_{02}x_2^2\)
specializes the residual.  For \(2^a\parallel N\), \(a\ge2\), the
rank-two coefficient \(A_{12}^{(N)}x_1x_2\) is not detected by cyclic
restrictions and is part of the finite-stabilizer orientation test of
Proposition~\ref{prop:finite-stabilizer-orientation-detection}.  The
orientation lift technology is
BBDJS \cite{BBDJS2015} on the unreduced side and Joyce--Upmeier
\cite{JoyceUpmeier,JoyceTanakaUpmeier} on the multiplicative side; the
reduced \(K3\times E\)-quotient theorem is separate.

\paragraph{(O1)\(^+\) — Weyl-equivariant determinant-line transport along
\(W^{(2)}(\La^{2,1}_{II})\) reflections.}
The chosen lifts
\(\tau_i:o_{\mathcal C}\to s_{\delta_i}^{*}o_{\mathcal C}\), \(i=1,2,3\),
of the three type-II wall transports satisfy
\([c_o]=0\) in
\(H^2(W^{(2)}(\La^{2,1}_{II});\mathbb F_2)\) (equivalently, in the
finite action-groupoid cohomology
\(H^2(\mathcal G_R;\underline{\mathbb F}_2)\) at every \(R\)),
\(\widetilde\tau_i^{\,2}=1\), and the quotient-orientation transport
defects \(\omega_{i,\mathcal C}=0\) for every simple reflection.  Then
\(\widetilde\tau_i\) define an honest Weyl action on the reduced quotient
orientation; the passage from the three generator signs to a character on
\(W^{(2)}(\La^{2,1}_{II})\) is the group-theoretic part of
Definition~\ref{def:O1-plus-weyl-transport}.

\paragraph{(O2) — Pfaffian local model and wall atlas on type-II walls.}
At a generic point of each
\(\Hpl_{\delta_i}\), the chosen real local equation
\(u_i=u_{\delta_i}\) with \(s_{\delta_i}(u_i)=-u_i\), the normal
determinant-complex splitting
\[
K^{\mathrm{nor}}_{\delta_i,R}\simeq
\bigl[\R\xrightarrow{u_i}\R\bigr],
\qquad
\operatorname{Pf}(K^{\mathrm{nor}}_{\delta_i,R})=\upsilon_{\delta_i,R}u_i,
\qquad
s_{\delta_i}(\upsilon_{\delta_i,R})=\upsilon_{\delta_i,R},
\]
and the rank-one normal-form rank \(N_{\delta_i}^{\mathrm{Pf}}=1\) hold.
The wall charts cover every retained component and boundary stratum of
\(\Hpl_{\delta_i}\); on overlaps the tangent/normal splittings, normal
Pfaffian units, quotient orientations, and protected integration maps
agree after a zero \v Cech sign class.  The retained atlas datum must
contain, for the middle wall \(\delta_2\leftrightarrow(0,1,1)\), a
wrapped wall object
\(x_{\delta_2,R}\in\mathcal M_{\eta,R}^{\mathrm{wr,rig}}\) with
\(b_R^{\mathrm{geom}}(\eta)>0\) and \(\overline\Pi_X(\eta)=(0,1,1)\).
Proposition~\ref{prop:appE-local-pfaffian-sign} supplies only the
rank-one local Pfaffian sign on such a chart; the component, boundary,
overlap, quotient, and protected-integration compatibilities are part of
\((O2_{\mathrm{atlas}})\).

\subsection{The OP minus sign is a scalar branch convention}
\label{subsec:ch6-op-scalar-branch}

The Oberdieck--Pixton (OP) signed scalar is
\(C_{\mathrm{OP},X}=-4096\).  Its absolute square-normalization is
\(C_{\square,X}=4096=64^2\), and its sign is the OP scalar branch
recorded by the fermionic shift \((-p_{\mathrm{DT}})^n\) of the
OP/DT chamber.  It is not the source of the orientation character
\(\epsilon_o\); orientation lines are killed by squaring.  The
orientation character is computed by monodromy
of the Pfaffian section around the type-II walls (Lemma
\ref{prop:appE-local-pfaffian-sign}); equivalently, by the
\((-1)^{N_{\delta_i}^{\mathrm{Pf}}}\) values from (O2).

\section{The Pfaffian--Dirac theorem}
\label{sec:ch6-pfaffian-dirac}

The protected Pfaffian of the Dirac--Igusa pro-object is compared with
the automorphic section.

\begin{theorem}[Pfaffian--Dirac section identity]
\label{thm:ch6-pfaffian-dirac}
Granted the four obstruction classes of
Section~\ref{subsec:ch6-obstructions}, with \((O2)\) understood as the
retained wall-atlas datum \((O2_{\mathrm{atlas}})\), and granted the
finite geometric Pfaffian datum \((P_{\mathrm{fin}})\) of
Definition~\ref{def:finite-geometric-pfaffian-datum} on a cofinal HN
system, including skew perfect complexes, Pfaffian line and section,
support, parity, active-support coverage, type-II wall divisor order,
leading coefficient, Pfaffian-line comparison, and Pfaffian-line
Mittag--Leffler compatibilities,
the protected Pfaffian of the Dirac--Igusa pro-object
\(\mathfrak D_X^{\mathrm{DI}}\) equals the automorphic section
\(\mathcal D_X=\Delta_5\):
\[
\boxed{\;
\iota_{\mathrm{aut}}\!\left(
\operatorname{Pf}_{\mathrm{prot}}(\mathfrak D_X^{\mathrm{DI}})\right)
\;=\;\mathcal D_X\;=\;\Delta_5
\;}
\]
as sections of \(\mathcal L^5\otimes\nu_{\Delta_5}\) over the genus-two
Siegel space \(\Sp(2,\Z)\backslash\mathcal H_2\), where
\[
\iota_{\mathrm{aut}}:\mathscr L_{\mathrm{Pf},X}
\xrightarrow{\,\sim\,}\mathcal L^5\otimes\nu_{\Delta_5}
\]
is the fixed Pfaffian-to-automorphic line comparison, compatible with
the cusp trivialization at \(\mathfrak c_\infty\) and with squaring to
\(\mathcal L^{10}\).  Equivalently, in the type-II Borcherds chamber,
\[
\iota_{\mathrm{aut}}\!\left(
\operatorname{Pf}_{\mathrm{prot}}(\mathfrak D_X^{\mathrm{DI}})\right)
\big|_{\mathfrak c_\infty}
=64\,q^{1/2}r^{1/2}s^{1/2}
\!\!\!\!\prod_{(n,l,m)\in\Gamma_{\mathrm{eff}}}\!\!\!
(1-q^nr^ls^m)^{f(nm,l)}.
\]
After squaring, the orientation forgets and the OP scalar branch
identifies the level-\(n\) protected scalar shadow with the inverse
squared Pfaffian:
\[
 Z^X_{\mathrm{OP}}\;=\;
 -\,64^2\,(\Pf_{\mathrm{prot}}(\mathfrak D_X^{\mathrm{DI}}))^{-2}
 \;=\;-4096\,\Delta_5^{-2},
\]
in agreement with \cite{OB}, where \(4096 = 64^2 =
([q^{1/2}r^{1/2}s^{1/2}]\Delta_5)^2\) is the squared theta-normalization
and the leading minus sign is the OP scalar branch convention.  The
BPS normalization rescales by \(-4096^{-1}\):
\(Z^{K3\times E}_{\mathrm{BPS}} = -4096^{-1}Z^X_{\mathrm{OP}}
= \Delta_5^{-2}\).
\end{theorem}

\begin{proof}
The inverse-limit assembly consists of a Pfaffian product, a leading
coefficient identification with a known weight-five Siegel form, and a
\(q\)-expansion principle on a punctured neighbourhood of the type-II
cusp \(\mathfrak c_\infty\).

\smallskip\emph{Finite-stage Pfaffian product
\((\delta+\varepsilon)\).}  The retained D0-HN datum supplies the
cosection-reduced d-critical pro-limit on the retained HN system.  The
finite geometric Pfaffian datum \((P_{\mathrm{fin}})\) supplies, at
every cofinal \(R\), the cosection-reduced skew perfect complexes
\(K^{\mathrm{red}}_{\mathcal C,R}\), the first-order block
\(\mathfrak D_{X,R}\) on the reduced compact quotient
\(\mathfrak Q_R=(\mathfrak M_R/E)^{\mathrm{red}}\), the Pfaffian line
\(\mathscr L_{\mathrm{Pf},R}\), the section \(\operatorname{pf}_{X,R}\),
active-support coverage, the support and parity residual vanishings,
the type-II wall divisor order, and the normalized leading coefficient
\(\mathfrak o^{\mathrm{lead}}_R=0\).  Hypothesis (O1) supplies
the orientation datum whose square-root line \(o_{\mathcal C}\) satisfies
\(o_{\mathcal C}^{\otimes 2}\simeq
\det\operatorname{RHom}_{\mathrm{red}}(\mathcal C,\mathcal C)\) and
extension-multiplicativity
\(o_{\mathcal C\oplus\mathcal C'}\simeq
o_{\mathcal C}\otimes o_{\mathcal C'}\), descended under the reduced
\(E\)-quotient.  The protected Pfaffian functor of (\(\beta\)) is then the
graded Pfaffian product of the determinant of
\(\operatorname{RHom}_{\mathrm{red}}\) over the retained
charge strata.  The exponent-matching clause of the retained
finite-stage Dirac block gives
\[
\sdim P^\Pi_{X,(n,l,m)}=f(nm,l)
\quad((n,l,m)\in\Gamma_{\mathrm{eff}}),
\]
so the graded Pfaffian product, expanded in the type-II cusp polarization
\((q,r,s)=(e^{2\pi iz_1},e^{2\pi iz_2},e^{2\pi iz_3})\), is
\[
\operatorname{Pf}_{\mathrm{prot}}(\mathfrak D_{X,R})
=v_0\!\!\prod_{(n,l,m)\in\Gamma_{\mathrm{eff,R}}}\!\!
(1-q^nr^ls^m)^{f(nm,l)},
\qquad v_0=64q^{1/2}r^{1/2}s^{1/2},
\]
where \(\Gamma_{\mathrm{eff,R}}\) is the active support cut by HN height
\(R\).  The Weyl monomial \(v_0\) is the theta normalization
\([q^{1/2}r^{1/2}s^{1/2}]\Delta_5=64\) of \cite[Theorem~4.1]{GN}.
The finite normalization table \(\texttt{delta5\_theta\_leading}\)
records the scalar equalities \(64^2=4096\), \(D_5=64^{-1}\Delta_5\),
and \(\chi_{10}^{\mathrm{OP}}=4096^{-1}\Delta_{10}\), together with
separate convention rows for
\(\Delta_{10}\), \(\chi_{10}\), \(\Phi_{10}^{\mathrm{un}}\), the OP
monic form, and the inverse scalar branches; it has no
Pfaffian-line, orientation-character, or O2 wall-atlas content.
This product is the cusp expansion of the supplied Pfaffian section, not
a replacement for that section: by
Proposition~\ref{prop:scalar-pfaffian-data-not-section}, the scalar
product, its square, and the OP scalar do not determine
\(\mathscr L_{\mathrm{Pf},R}\), \(\operatorname{pf}_{X,R}\), or
\(\epsilon_{o,R}\) without the Pfaffian-line and line-comparison data.

\smallskip\emph{Inverse limit \((\gamma)\).}  The
pro-object ambient is the cofinal HN system; the retained D0-HN datum
and finite geometric Pfaffian datum supply strict transition maps with
vanishing successor residuals
\(\mathfrak o^{\mathrm{atlas}}_{R^+R}=
\mathfrak o^{\mathrm{tr}}_{R^+R}=0\); the finite Pfaffian sections
\(\operatorname{pf}_{X,R}\) form a Mittag-Leffler system on the active
support, and the inverse limit is the formal Borcherds product
\[
\operatorname{Pf}_{\mathrm{prot}}(\mathfrak D_X^{\mathrm{DI}})
\big|_{\mathfrak c_\infty}
=\lim_R\operatorname{Pf}_{\mathrm{prot}}(\mathfrak D_{X,R})
=64\,q^{1/2}r^{1/2}s^{1/2}
\!\!\!\prod_{(n,l,m)\in\Gamma_{\mathrm{eff}}}\!\!
(1-q^nr^ls^m)^{f(nm,l)}.
\]
Class \(M\) completion (chain-level after weight-completion, in the
sense of Vol I's class \(M\) global triangle) is exactly the ambient in
which \(\lim_R\) commutes with the graded Pfaffian product on the
active support; this is the \(\gamma\) clause.

\smallskip\emph{Comparison with the imported Borcherds product
\((\beta)\).}  By
Proposition~\ref{prop:bgn-identification-recap},
the formal product
\(64\,q^{1/2}r^{1/2}s^{1/2}\prod(1-q^nr^ls^m)^{f(nm,l)}\) equals the
Fourier expansion of the automorphic section
\(\mathcal D_X=\Delta_5\) at \(\mathfrak c_\infty\); this is the
Borcherds--Gritsenko--Nikulin theorem
\cite{Borcherds95,GN,GNII}.  The fixed
\(\iota_{\mathrm{aut}}:\mathscr L_{\mathrm{Pf},X}\xrightarrow{\sim}
\mathcal L^5\otimes\nu_{\Delta_5}\) of (\(\beta\)), compatible with cusp
trivialization and squaring, transports the inverse-limit Pfaffian to a
section of \(\mathcal L^5\otimes\nu_{\Delta_5}\).

\smallskip\emph{\(q\)-expansion principle \((\alpha)\).}
The chart \(\alpha\) is the open factorization category
\(\mathcal C\) on \((X,D,\tau)\) with boundary vacuum \(b\): on the
punctured neighbourhood of \(\mathfrak c_\infty\) inside Sp\((2,\Z)
\backslash\mathcal H_2\) the chart algebra
\(A_b=\mathrm{End}_{\mathcal C}(b)\) supplies the local
factorization-algebra structure within which the Pfaffian section is
holomorphic.  Two holomorphic genus-two Siegel modular forms with the
same Maass character \(\nu_{\Delta_5}\) and the same Fourier expansion at
\(\mathfrak c_\infty\) are equal: this is the \(q\)-expansion principle
for holomorphic Siegel modular forms
\cite{Igusa1962,Igusa1964II}.  Both sides
have weight \(5\) and Maass character \(\nu_{\Delta_5}\): the right-hand
side by the Borcherds--Gritsenko--Nikulin theorem; the left-hand side
because the orientation character of the retained
\(\mathfrak D_X^{\mathrm{DI}}\), under (O1)+(O1)\(^+\)+(O2), is exactly
\(\nu_{\Delta_5}\) on \(W^{(2)}(\La^{2,1}_{II})\) (this is
Theorem~\ref{thm:ch6-orientation-character}).  Hence
\(\iota_{\mathrm{aut}}(\operatorname{Pf}_{\mathrm{prot}}(\mathfrak D_X^{\mathrm{DI}}))=\Delta_5\)
as sections of \(\mathcal L^5\otimes\nu_{\Delta_5}\).

\smallskip\emph{Squared scalar \((\beta)\).}
Squaring forgets orientation: the squared line
\(\bigl(\mathcal L^5\otimes\nu_{\Delta_5}\bigr)^{\otimes 2}=\mathcal L^{10}\)
is trivially charactered, and \(\Delta_5^2=\Delta_{10}\) is scalar.
The OP scalar normalization of Convention~\ref{conv:op-chamber} fixes the partition function
\(Z^X_{\mathrm{OP}}=-(\chi_{10}^{\mathrm{OP}})^{-1}=-4096\,\Delta_5^{-2}\),
with \(\chi_{10}^{\mathrm{OP}}=D_5^2=4096^{-1}\Delta_5^2\) and
\(D_5=64^{-1}\Delta_5\).  Composition of the squared Pfaffian section
with the OP scalar branch gives the OP chamber scalar representative
\[
 Z^X_{\mathrm{OP}}\;=\;-\,(\Pf_{\mathrm{prot}}(\mathfrak D_X^{\mathrm{DI}}))^{-2}\cdot
 C_{\square,X},\qquad C_{\square,X}=4096,
\]
identifying \(Z^X_{\mathrm{OP}}=-4096\,\Delta_5^{-2}\); the minus sign
is the OP scalar branch convention, recorded by
\((-p_{\mathrm{DT}})^n\), and is not a feature of the orientation
character.  The protected scalar trace is
\(Z_{\mathrm{BPS}}^{K3\times E}=\Delta_5^{-2}\).
\end{proof}

\begin{remark}[Why (D0)+(O1)+(O1)\(^+\)+(O2)+\((P_{\mathrm{fin}})\)]
\label{rem:ch6-why-four}
Each named datum is essential.  (D0) gives the existence and
the finite first-order block; without it there is no
\(\mathfrak D_{X,R}\) and no Pfaffian product to take.  (O1) gives
the orientation line and the reduced \(E\)-quotient descent; without
the line the Pfaffian section has no domain.  (O1)\(^+\) gives the
honest Weyl action on the reduced quotient orientation; without it the
type-II reflection signs are projective and the orientation character is
projective, not a homomorphism.  (O2) gives the local Pfaffian normal
form; without it the rank \(N_{\delta_i}^{\mathrm{Pf}}\) is undetermined
and the local sign cannot be computed.  The finite geometric Pfaffian
datum \((P_{\mathrm{fin}})\) gives the skew perfect complexes,
Pfaffian line, section, support, parity, active-support coverage, wall
divisor order, leading coefficient, line comparison, and
Pfaffian-line Mittag--Leffler compatibility; without it the inverse-limit
product has no section-level domain and no determinant-line object.  If
one of these data is omitted, one of the proof steps loses its object.
\end{remark}

\section{The orientation character identity}
\label{sec:ch6-orientation-character}

The second theorem identifies the geometric orientation monodromy character
\(\epsilon_o\) of the Dirac--Igusa pro-object with the automorphic Maass
character \(\nu_{\Delta_5}\) on the type-II Weyl group.  The proof
splits into a local Pfaffian wall computation, a group-theoretic
extension, and an independent automorphic check.

\begin{theorem}[Orientation character]
\label{thm:ch6-orientation-character}
Recall the type-II simple roots
\(\delta_1=2f_2-f_3,\delta_2=2f_{-2}-f_3,\delta_3=f_3\) on
\(\La^{2,1}_{II}\subset\La^{2,1}=\Z f_2\oplus\Z f_3\oplus\Z f_{-2}\)
(Equation~\eqref{eq:bkm-type-ii-simple-roots}).  Granted the orientation
classes (O1), (O1)\(^+\), (O2), the geometric orientation monodromy
character of the Dirac--Igusa pro-object equals the automorphic Maass
character on the type-II Weyl group:
\[
\boxed{\;
\epsilon_o\;=\;\nu_{\Delta_5}
\quad\text{on}\quad W^{(2)}(\La^{2,1}_{II})\;}.
\]
The character takes the value \(-1\) on each of the three type-II simple
reflections \(s_{\delta_1},s_{\delta_2},s_{\delta_3}\), and the
identity is reduced to these three values by the abelianization
\(W^{(2)}(\La^{2,1}_{II})_{\mathrm{ab}}\simeq(\Z/2)^3\)
(Definition~\ref{def:O1-plus-weyl-transport}).
\end{theorem}

The identity is a comparison of two characters defined independently:
\(\epsilon_o\) is computed by transport of the Pfaffian section under
the chosen Weyl action of (O1)\(^+\); \(\nu_{\Delta_5}\) is the Maass
multiplier of the holomorphic Siegel form \(\mathcal D_X=\Delta_5\) computed from
the multiplier system of the corresponding theta lift
\cite{MAASS:1,GN,Gritsenko99}.  No scalar branch convention enters this
sign comparison: replacing \(\Delta_5\) by \(c\Delta_5\),
\(c\in\C^\times\), changes neither the Maass character nor the divisor
orders.  The factor \(64\) names the monic Borcherds product
\(D_5=64^{-1}\Delta_5\); the OP/DT scalar \(-4096\) is unused.

\begin{proof}
Three inputs combine:
Proposition~\ref{prop:appE-local-pfaffian-sign} computes the local Pfaffian
signs on the three simple type-II walls;
Definition~\ref{def:O1-plus-weyl-transport} extends a generator
sign-tuple \((\epsilon_1,\epsilon_2,\epsilon_3)\in\{\pm1\}^3\) to a
unique character on \(W^{(2)}(\La^{2,1}_{II})\); and
Lemma~\ref{lem:bkm-maass-restriction} computes the automorphic
target values \(\nu_{\Delta_5}(s_{\delta_i})=-1\) directly from the
Maass multiplier and the Borcherds divisor order
\(d_{\delta_i}^{\mathrm{aut}}=1\).

\smallskip\emph{Local Pfaffian signs from \((O2)\).}  By (O2), at a
generic point of each \(\Hpl_{\delta_i}\), \(i=1,2,3\), the chosen
local equation \(u_i\) satisfies \(s_{\delta_i}(u_i)=-u_i\) and the
normal determinant-complex splits as
\(K^{\mathrm{nor}}_{\delta_i}\simeq[\R\xrightarrow{u_i}\R]\) with
\(\operatorname{Pf}=\upsilon_{\delta_i}u_i\) and
\(s_{\delta_i}(\upsilon_{\delta_i})=\upsilon_{\delta_i}\).  Transport
along the closed path
\(\ell_{\delta_i}\) (the path in the reduced quotient wall-complement
carrying \(u_i=\epsilon\) to \(u_i=-\epsilon\), closed by the (O1)\(^+\)
lift \(\widetilde\tau_i\)) sends
\(\upsilon_{\delta_i}u_i\) to \(-\upsilon_{\delta_i}u_i\), giving
\[
\epsilon_o(s_{\delta_i})=(-1)^{N_{\delta_i}^{\mathrm{Pf}}}=-1,
\qquad i=1,2,3.
\]
This is Proposition~\ref{prop:appE-local-pfaffian-sign}.  Squaring eliminates
the sign and hence eliminates the OP scalar branch ambiguity: the squared
Pfaffian \(\operatorname{Pf}(K^{\mathrm{nor}}_{\delta_i})^2=u_i^2\) is
fixed by \(s_{\delta_i}\).

\smallskip\emph{Group-theoretic extension via \((O1)^+\).}  The
type-II Weyl group has Coxeter presentation
\[
W^{(2)}(\La^{2,1}_{II})
\simeq
\langle s_{\delta_1},s_{\delta_2},s_{\delta_3}\mid
s_{\delta_1}^2=s_{\delta_2}^2=s_{\delta_3}^2=1\rangle,
\]
with no braid relations — the off-diagonal entries
\((\delta_i,\delta_j)=-2\) for \(i\ne j\) give Coxeter exponents
\(m_{ij}=\infty\).  Hence
\(W^{(2)}(\La^{2,1}_{II})_{\mathrm{ab}}\simeq(\Z/2)^3\) with one
generator per \(s_{\delta_i}\), and any \(\{\pm1\}\)-valued assignment
on the generators extends uniquely to a homomorphism.  By (O1)\(^+\),
\([c_o]=0\), \(\widetilde\tau_i^{\,2}=1\), and
\(\omega_{i,\mathcal C}=0\); the lifts \(\widetilde\tau_i\) define an
honest action on the reduced quotient orientation, and the induced
character has \(\epsilon_o(s_{\delta_i})=-1\) for \(i=1,2,3\).  This
is Definition~\ref{def:O1-plus-weyl-transport}.

\smallskip\emph{Independent automorphic computation of
\(\nu_{\Delta_5}\).}  The Maass multiplier
\(\nu_{\Delta_5}\) of the singular theta lift of
\(\phi^{\mathrm{Weil}}_{0,1}\) is
computed from the multiplier system of the Borcherds product (Borcherds
\cite{Borcherds95}, Gritsenko--Nikulin
\cite[Section~3]{GN},\cite{Gritsenko99}, Maass \cite{MAASS:1}).
On the type-II simple roots
\(\delta_1=2f_2-f_3\), \(\delta_2=2f_{-2}-f_3\), \(\delta_3=f_3\), the
Borcherds divisor orders are
\[
d_{\delta_i}^{\mathrm{aut}}
:=\ord_{\Hpl_{\delta_i}}(\Delta_5)=f(0,\pm1)=1,
\qquad i=1,2,3,
\]
so \(\nu_{\Delta_5}(s_{\delta_i})=(-1)^{d_{\delta_i}^{\mathrm{aut}}}=-1\)
for \(i=1,2,3\).  This is
Lemma~\ref{lem:bkm-maass-restriction}.  The chamber-permuting
roots \(f_{-2}-f_2,f_2-f_3,f_{-2}-f_3\) lie in
\(\La^{2,1}\setminus\La^{2,1}_{II}\) and are not type-II walls; their
Maass values are \(+1\) but no Hall orientation sign is asserted on them
unless a semidirect-product Hall lift is constructed
(Remark~\ref{rmk:ch6-S3-extension}).

\smallskip\emph{The comparison \(\epsilon_o=\nu_{\Delta_5}\).}
Two \(\{\pm1\}\)-valued characters of \(W^{(2)}(\La^{2,1}_{II})\) that
agree on \(s_{\delta_1},s_{\delta_2},s_{\delta_3}\) agree everywhere
by the Coxeter presentation above.  The local Pfaffian computation and
the automorphic computation give that both \(\epsilon_o\) and
\(\nu_{\Delta_5}\) take the value \(-1\) on each of the three simple
reflections.  Therefore
\(\epsilon_o=\nu_{\Delta_5}\) on
\(W^{(2)}(\La^{2,1}_{II})\).
\end{proof}

\begin{remark}[The OP minus sign is not the source of \(\epsilon_o\)]
\label{rmk:ch6-op-not-orientation-source}
The local Pfaffian computation and the automorphic multiplier
computation are independent.  The former uses no automorphic data; the
latter uses no Pfaffian local model.  In particular,
the OP scalar
\(C_{\mathrm{OP},X}=-4096=(-1)\cdot C_{\square,X}\) does
not source the orientation character: squaring forgets orientation, so
\(C_{\mathrm{OP},X}\) carries no information about \(\epsilon_o\).  The
minus sign \(\operatorname{sign}(C_{\mathrm{OP},X})=-1\) is the OP scalar
branch convention of \cite{OB}, recorded by the fermionic
shift \((-p_{\mathrm{DT}})^n\) of the OP/DT chamber.  The orientation
character is a \emph{geometric} monodromy datum, computed by transport
of the Pfaffian line around type-II walls; it is fixed by the data of
(O1)+(O1)\(^+\)+(O2), and is independently checked against the
automorphic target by the Maass multiplier computation.
\end{remark}

\begin{remark}[Imaginary roots are not Weyl reflections]
\label{rmk:ch6-imaginary-roots}
The Weyl group \(W^{(2)}(\La^{2,1}_{II})\) of the type-II BKM datum is
generated by real simple reflections only.  The imaginary simple roots
contribute to the Weyl--Kac--Borcherds denominator through the
multiplicity coefficients \(m(a)\) and \(\tau(a)\) of
Gritsenko--Nikulin \cite[Theorem~2.1]{GNII},\cite[Section~3]{GN}; they
do not define reflections in
\(W^{(2)}(\La^{2,1}_{II})\).  Hence
\(\epsilon_o:W^{(2)}(\La^{2,1}_{II})\to\{\pm 1\}\) is a real-root
reflection datum.  Local divisor monodromy of a compact Hall theory
around imaginary components of the Borcherds divisor is not a Weyl
character and is not asserted here.
\end{remark}

\begin{remark}[Extension to \(W^{(2)}(\La^{2,1}_{II})\rtimes\Aut(\Poly_{II})\)]
\label{rmk:ch6-S3-extension}
The chamber-permuting symmetry group
\(\Aut(\Poly_{II})\simeq S_3\) acts on the Coxeter generators
\(s_{\delta_1},s_{\delta_2},s_{\delta_3}\) by permutation; its three
reflections \(s_{f_{-2}-f_2},s_{f_2-f_3},s_{f_{-2}-f_3}\) have
\(\nu_{\Delta_5}=+1\).  An extension of \(\epsilon_o\) to
\(W^{(2)}(\La^{2,1}_{II})\rtimes\Aut(\Poly_{II})\) would require a
separate Hall semidirect-product lift of the chamber-permuting
reflections.
\end{remark}

\section{The primitive recognition theorem}
\label{sec:ch6-primitive-recognition}

The third theorem is a criterion for identifying the source-side
primitive Lie superalgebra of the Dirac--Igusa pro-object with the BKM
denominator algebra \(\mathfrak g_{\Delta_5}\).  The conditions form the
\emph{recognition criterion}: ten finite-stage data carrying compact
geometric provenance, none of which is inferable from the Borcherds
product or signed determinant alone.

\begin{theorem}[Primitive recognition]
\label{thm:ch6-primitive-recognition}
Granted the recognition criterion (S1)--(S10) of
Theorem~\ref{thm:primitive-recognition}, the cofinal finite-window
datum of Definition~\ref{def:cofinal-primitive-recognition-datum},
the chiral Koszul source datum of
Definition~\ref{def:chiral-koszul-source-datum}, the finite Koszul
comparison datum of
Definition~\ref{def:finite-koszul-comparison-datum}, and the retained
orientation data \((O1_{\mathrm{quot}})\), the chosen Weyl lifts, and
\((O2_{\mathrm{atlas}})\),
suppose the compact \(K3\times E\) Hall/factorization construction
supplies, on a cofinal sequence of downward-saturated finite root windows
\(W_1\subset W_2\subset\cdots\) with
\(\bigcup_\nu W_\nu=\mathcal R_+\):
\begin{itemize}
\item \emph{(S1) Chain-level descent endpoint.}  The (D0)-degenerated
  source coalgebra \((C_{X,W_\nu},Q_{X,W_\nu})\) of
  Definition~\ref{def:chiral-koszul-source-datum}, descended chain-level
  to its primitive part on every finite window
  \(W_\nu\), carrying both the orientation lane and the
  Frobenius--Serre primitive structure feeding the recognition window.
\item \emph{(S2) Cartan and root-degree matching.}  An even Cartan
  identification \(\iota_H:P^{\Pi,\mathrm{red}}_{X,0}\xrightarrow{\sim}
  \La^{2,1}_{II}\otimes\C\), with Gram-grading
  \(\alpha(n,l,m)=2nf_2-lf_3+2mf_{-2}\).
\item \emph{(S3) Simple primitive representatives.}  Parity-homogeneous
  \(e_i^X,f_i^X,h_i^X\) for every simple-root index \(i\in I\) of
  \(\mathscr B_{\Delta_5}\), with Gritsenko--Nikulin parity dimensions
  \(\mu_{\overline 0}(a),\mu_{\overline 1}(a)\) on imaginary fibres.
\item \emph{(S4) Borcherds--Kac relations.}  Cartan, Chevalley
  (\([e_i^X,f_j^X]=\delta_{ij}h_i^X\)), real Serre, Borcherds
  orthogonality, super sign, and isotropic-string relations satisfied in
  the protected Hall bracket.
\item \emph{(S5) Completed Hopf radical equality.}  The completed Hopf
  pairing has closed homogeneous radical \(\operatorname{Rad}_X\) which
  is a Lie ideal and a coideal, and identifies under the radical
  quotient with \(\operatorname{Rad}_{\mathrm{GN}}\).
\item \emph{(S6) No-extra-relations.}  The presentation map
  \(\pi:\widehat{\mathfrak F}(\mathscr B_{\Delta_5})\to
  P_X^{\Pi,\mathrm{red}}\) has
  \(\ker\pi=\overline{J_{\mathrm{BK}}+\operatorname{Rad}_{\mathrm{GN}}}\),
  on every finite window separately.
\item \emph{(S7) Generation.}  \(P_X^{\Pi,\mathrm{red}}\) is
  topologically generated by \(P^{\Pi,\mathrm{red}}_{X,0}\) and
  \(\{e_i^X,f_i^X\}_{i\in I}\).
\item \emph{(S8) Completed PBW.}  The associated graded
  \(\operatorname{gr}_{\mathrm{PBW}}U(P_X^{\Pi,\mathrm{red}})_W
  \xrightarrow{\sim}
  \operatorname{gr}_{\mathrm{PBW}}U(\mathfrak g_{\Delta_5})_W\) is an
  isomorphism on every finite window, with strict transitions.
\item \emph{(S9) Exact triangular completion.}  The HN inverse limit is
  exact on the saturated finite root windows: \(\lim^1\) vanishes for
  the kernels, radicals, PBW associated gradeds, and root-space parity
  pieces.
\item \emph{(S10) Full parity dimensions.}  Two integers per nonzero
  root degree, in both signs:
  \(\dim P^{\Pi,\mathrm{red}}_{X,\gamma,\overline p}
  =\rootmult_{\overline p}^{\mathrm{GN}}(\alpha(\gamma))\),
  \(p=0,1\), and the same equality in degree \(-\gamma\) through the
  Borcherds--Kac Chevalley anti-involution; not only their signed
  difference.
\end{itemize}
Then the recognition datum determines an isomorphism of completed
\(\Gamma_{\mathrm{gram}}\)-graded Lie superalgebras
\[
\boxed{\;
H^0\Prim_{\mathrm{prot}}\!\bigl(
\overline\Pi_{X,*}^{\Theta}\mathcal F_X\bigr)
\big/\overline{\operatorname{Rad}\langle\,\cdot\,,\,\cdot\,\rangle_X}
\;\cong\;\mathfrak g_{\Delta_5}
\;}
\]
and the Weyl--Kac--Borcherds denominator of
\(P_X^{\Pi,\mathrm{red}}\) is \(\operatorname{den}(\mathfrak g_{\Delta_5})
=64^{-1}\Delta_5(2Z)\).  If in addition the chiral Koszul source datum of
Definition~\ref{def:chiral-koszul-source-datum} is fixed, including
the finite Koszul comparison datum of
Definition~\ref{def:finite-koszul-comparison-datum} and the strict
Koszul Mittag--Leffler exactness hypothesis of
Proposition~\ref{prop:chiral-koszul-source-datum-consequence}, then the
chiral cobar of the source coalgebra \(C_X\) is identified with the
formal current envelope compatibly with the Weyl action, Pfaffian
orientation character, Hall pairing, radical quotient, and PBW
associated graded:
\[
\Omega_E^{\mathrm{ch}}C_X\;\simeq\;\mathsf{Den}_{\Delta_5,E}.
\]
\end{theorem}

\begin{proof}
The argument is the inverse-limit assembly of (S1)--(S10) on each
finite window, together with the chiral cobar identification.  The
window-wise theorem is Theorem~\ref{thm:primitive-recognition}, with
the algebraic finite-window step isolated in
Proposition~\ref{prop:finite-window-presentation-recognition}; the
additional input is the chiral-cobar comparison
\(\Omega_E^{\mathrm{ch}}C_X\simeq\mathsf{Den}_{\Delta_5,E}\), with
the compatibility statement supplied by
Proposition~\ref{prop:finite-koszul-comparison-compatibility} and its
Mittag--Leffler passage in
Proposition~\ref{prop:chiral-koszul-source-datum-consequence}.

\smallskip\emph{Presentation map and kernel
\((\alpha+\beta)\).}  (S2) and (S3) define the continuous Lie-superalgebra
homomorphism
\[
\pi:\widehat{\mathfrak F}(\mathscr B_{\Delta_5})
\longrightarrow P_X^{\Pi,\mathrm{red}}
\]
on the descended Gram-graded completion.  (S4) puts the Borcherds--Kac
ideal \(J_{\mathrm{BK}}\) in the kernel.  (S6) gives the kernel equality
\(\ker\pi=\overline{J_{\mathrm{BK}}+\operatorname{Rad}_{\mathrm{GN}}}\)
on every finite window, hence injectivity of the descended map
\(G(\mathscr B_{\Delta_5})\to P_X^{\Pi,\mathrm{red}}\).  (S7) gives
surjectivity.  Thus the descended map is an isomorphism on every finite
window by
Proposition~\ref{prop:finite-window-presentation-recognition}.

\smallskip\emph{Radical descent \((\varepsilon)\).}
(S5) identifies the source Hopf radical
\(\overline{\operatorname{Rad}_X}\) with the target radical
\(\overline{\operatorname{Rad}_{\mathrm{GN}}}\).  The required Hopf
adjointness, ideal/coideal property, quotient-tensor non-degeneracy, and
closedness on the inverse system are part of (S5), equivalently the
finite matrix conditions (R4)--(R5).  The orientation lane supplies the
trace pairing and Frobenius formalism on which those equations are
written; it does not prove the source matrices \(M,D,B,G,K,Q\), the
kernel equality, or the PBW comparison.

\smallskip\emph{Full parity dimensions \((\delta)\).}
(S10) is the two-integer equality
\(\dim P^{\Pi,\mathrm{red}}_{X,\gamma,\overline p}=
\rootmult_{\overline p}^{\mathrm{GN}}(\alpha(\gamma))\) for \(p=0,1\), not
just the signed difference.  This rules out hidden zero-superdimension
source summands, wrong parity lifts invisible to the determinant, and
parity-cancelling ideals.  A putative compact source matching every
Hall bracket template of \(\mathfrak g_{\Delta_5}\) on a finite window
can still fail (S5): a kernel element \(k\in K_{\gamma_\rho}\) with
\((Q_{\gamma_\mu}\otimes Q_{\gamma_\nu})D^{\mu,\nu}_\rho(k)\ne 0\) is
invisible to the quotient target table yet breaks the radical
coideal property, so the Hopf quotient is not defined
(Theorem~\ref{thm:primitive-recognition}).  This is the finite
tensor-Hopf obstruction hidden by signed determinant data; (S5)+(S10)
together block it.

\smallskip\emph{PBW comparison and exact completion
\((\gamma)\).}  (S8) is the associated-graded isomorphism on every
finite window, with strict transitions.  (S9) is the exact-completion
clause: \(\lim^1\) vanishes for kernels, radicals, PBW associated
gradeds, and root-space parity pieces.  Equivalently, no new generator,
relation, radical element, or parity-cancelling summand appears only
after completion.  Proposition~\ref{prop:strict-ml-completion-primitive-recognition}
is the homological algebra that passes from the finite-window
isomorphisms to the completed source and target.  Class \(M\)
completion is the ambient: chain-level identifications hold after
weight-completion; this is the \(\gamma\) clause.

\smallskip\emph{Passage to the inverse limit.}  Since the
windows are cofinal (\(\bigcup_\nu W_\nu=\mathcal R_+\)) and the radical
in (S5) is the closure of the finite kernels, the finite isomorphisms
of the finite-window clauses determine the inverse-limit isomorphism
\[
H^0\Prim_{\mathrm{prot}}\!\bigl(
\overline\Pi_{X,*}^{\Theta}\mathcal F_X\bigr)
\big/\overline{\operatorname{Rad}\langle\,\cdot\,,\,\cdot\,\rangle_X}
=P_X^{\Pi,\mathrm{red}}
\xrightarrow{\,\sim\,}\mathfrak g_{\Delta_5}.
\]
The Weyl--Kac--Borcherds denominator of
\(P_X^{\Pi,\mathrm{red}}\) is therefore
\(\operatorname{den}(\mathfrak g_{\Delta_5})=64^{-1}\Delta_5(2Z)\), by
Theorem~\ref{thm:bkm-denominator-algebra-identity}.

\smallskip\emph{Chiral cobar identification.}  Suppose the
chiral Koszul source datum of
Definition~\ref{def:chiral-koszul-source-datum} is fixed:
\((\mathcal F_{X,\sigma,S,\le R}^{\mathrm{hyb}},C_{X,R},
F^{\mathrm{co}}_{R,\bullet},\Delta_R^{\mathrm{ch}},
b_{X,R},\Theta_{\mathrm{Kos},R})\) with vanishing residual
\(\mathfrak O_{\mathrm{Kos},R}=0\), together with the finite Koszul
comparison datum
\(\mathfrak K^{\mathrm{cmp}}_{X,R}\) and
\(\mathfrak O_{\mathrm{Kcmp},R}=0\), compatibly in \(R\) and with the
strict Mittag--Leffler exactness hypothesis of
Proposition~\ref{prop:chiral-koszul-source-datum-consequence}.  The
target side is the formal current envelope
\(\mathsf{Den}_{\Delta_5,E}=U_E^{\mathrm{ch}}\!
\bigl(\operatorname{Cur}_E(\mathcal A_{\mathrm{den}})\bigr)\)
of Proposition~\ref{prop:bkm-formal-current-envelope}.  The cobar comparison
of Propositions~\ref{prop:finite-koszul-comparison-compatibility}
and~\ref{prop:chiral-koszul-source-datum-consequence}
identifies
\[
\Omega_E^{\mathrm{ch}}C_X\;\simeq\;\mathsf{Den}_{\Delta_5,E},
\]
through the source-internal bar/cobar admissibility (the residual
\(\mathfrak o^{\varepsilon,\mathrm{src}}_R=0\)) and the source-to-target
quasi-isomorphism (\(\mathfrak o^{\vartheta}_R=0\)), and it preserves
the Weyl, Pfaffian-orientation, Hall-pairing, radical-quotient, and
PBW structures used above.  Agreement on the
primitive summand is not enough by
Proposition~\ref{prop:primitive-restriction-not-koszul}; the full cone
of \(\vartheta_R\) and all compatibility defect complexes must be
acyclic.  The construction avoids the
target-internal bar/cobar counit; this is the \(\beta\)-separation
between observables and states.
\end{proof}

\begin{remark}[The zero-bracket graded super-vector-space counterexample]
\label{rmk:ch6-zero-bracket-counterexample}
The criterion (S1)--(S10) blocks the trivial counterexample of a graded
super vector space \(V_\bullet\) with all brackets zero and parity
dimensions matching \(\rootmult_{\overline p}^{\mathrm{GN}}\) on every
window.  Such a \(V_\bullet\) may satisfy the target parity table and
the signed-superdimension residual on the active support.  It fails
(S4): the Chevalley relation
\([e_i^X,f_j^X]=\delta_{ij}h_i^X\) fails on the diagonal.  It fails
(S6): the kernel of the presentation map
\(\pi:\widehat{\mathfrak F}(\mathscr B_{\Delta_5})\to V_\bullet\) is not
\(\overline{J_{\mathrm{BK}}+\operatorname{Rad}_{\mathrm{GN}}}\) but
contains every commutator \([e_i,e_j]\) with \((\alpha_i,\alpha_j)\ne 0\),
so the induced map \(G(\mathscr B_{\Delta_5})\to V_\bullet\) is not
injective.  It fails (S5) for the zero pairing because the radical is
all of \(V_\bullet\), so the radical quotient is zero and not the BKM
target.  If one installs a nonzero pairing by hand, the denominator
still does not supply its invariant Hopf adjointness, coideal radical,
or equality with \(\operatorname{Rad}_{\mathrm{GN}}\).  The
counterexample shows that signed-superdimension matching alone is not
recognition; the bracket-level data of (S4)+(S5)+(S6) is the
mathematical content.
\end{remark}

\begin{remark}[Recognition without (D0)+(O1)+(O1)\(^+\)+(O2)]
\label{rmk:ch6-recognition-without-orientation}
The recognition criterion of Theorem~\ref{thm:ch6-primitive-recognition}
nominally lists no orientation hypothesis among (S1)--(S10).  The
orientation classes enter through (S1) — the chain-level descent
\(\overline\Pi_{X,*}^{\Theta}\mathcal F_X\) is conditional on the
orientation lane via Definition~\ref{def:O1-strong-reduced-orientation} and the
Frobenius identity from the CY-3 Serre functor — and through (S5),
where the Hopf radical's coideal property is fixed by the same Frobenius
identity (Theorem~\ref{thm:hall-product}).  In the
chapter's combined statement
(Theorem~\ref{thm:ch6-pfaffian-dirac}+Theorem~\ref{thm:ch6-orientation-character}+Theorem~\ref{thm:ch6-primitive-recognition})
the orientation classes enter all three theorems.
\end{remark}

\section{Compatibility with chiral Koszul duality and Calabi--Yau quantum groups}
\label{sec:ch6-chiral-compatibility}

The three theorems sit inside chiral Koszul duality and
Calabi--Yau quantum groups as the View~\textcircled{2} /
View~\textcircled{3} face of the Calabi--Yau-to-chiral comparison.  Three
compatibilities are required.

\subsection{Primitive-to-shadow comparison}
\label{subsec:ch6-primitive-to-shadow}

The primitive-to-shadow comparison is the cross-level identity
\textcircled{\scriptsize 2}\(\leftrightarrow\)\textcircled{\scriptsize 3}
applied to the primitive-to-shadow part of the chain.  The Vol II primitive
datum on \((X,D,\tau)\) is the bicoloured nine-tuple
\[
\bigl(\mathcal C\big|_{(X,D,\tau)},\,
\mathcal D^{\mathrm{cl}}\big|_X;\,
b,\,A_b,\,F^{\mathrm{cl}};\,
\sigma_{\mathrm{half}},\,\operatorname{tr}^{\mathrm{cl}}_X;\,
\operatorname{Tr}_{\mathcal C}\bigr).
\]
The chart of vacuum \(b\) is the \(\alpha\)-datum; the chart
algebra \(A_b=\mathrm{End}_{\mathcal C}(b)\) is the open primitive
algebra; the closed-colour vacuum factorization algebra
\(F^{\mathrm{cl}}\) gives the closed-side trace
\(\operatorname{tr}^{\mathrm{cl}}_X\) which evaluates at the type-II
cusp \(\mathfrak c_\infty\) to the weight-\(5\) Borcherds--Gritsenko--Nikulin
section \(\mathcal D_X=\Delta_5\) (View \textcircled{\scriptsize 1}).  The
comparison
\(A_b=\Phi^{(\Sigma_{d-1},C)}_d(\mathcal C_X)\) is the chain-fusion datum of
chiral Koszul duality and Calabi--Yau quantum groups at level
\(2\): under that datum the Calabi--Yau-side chiral algebra
\(A_X\) is compared with the open-side chart algebra \(A_b\) on the same
curve \(C\) with boundary vacuum induced by
\((\mathcal C_X,\Sigma_{d-1},C)\).  In the present setting, the
comparison concerns the chiral algebra on \(E\) attached to the Hall
source of \(K3\times E\); conditional on the chiral Koszul source
datum, the finite Koszul comparison datum, and strict Koszul
Mittag--Leffler exactness, the chiral cobar
\(\Omega_E^{\mathrm{ch}}C_X\simeq\mathsf{Den}_{\Delta_5,E}\) of
Theorem~\ref{thm:ch6-primitive-recognition} is the source-side chiral
expression at level \(2\), and its constant-current Lie superalgebra is
\(\mathfrak g_{\Delta_5}\) (View \textcircled{\scriptsize 2}).  The
protected Pfaffian \(\operatorname{Pf}_{\mathrm{prot}}(\mathfrak D_X^{\mathrm{DI}})\)
is the pre-trace Pfaffian section.  After squaring and applying the
retained level-\(\mathsf Z\) datum, Vol~I Theorem~H
(chiral-Hochschild concentration on the Koszul locus) and Vol~II's
cyclic-Hochschild Beilinson stratification \textnormal{(i)+(ii)}
identify the derived-centre trace pairing of degree
\(-\dim_\C(K3\times E)\) whose scalar is \(\Delta_5^{-2}\)
(Appendix~\ref{app:obstruction-discharges},
Theorem~\ref{thm:G-vol2-discharge}).  The level-\(\mathsf A\)
gravity-line operator algebra acting on the boundary is the open
Hall--Borcherds residual of Vol~II and is not constructed by the
level-\(\mathsf Z\) trace criterion.

\subsection{Vol III chain fusion at level \(2\)}
\label{subsec:ch6-vol3-chain-fusion}

The Vol III \(\kappa\)-stratification supplies the
\(\kappa\)-tuple
\((\kappa_{\mathrm{cat}},\kappa_{\mathrm{ch}}^{\mathrm{Hodge}},
\kappa_{\mathrm{ch}}^{\mathrm{Heis}},\kappa_{\mathrm{BKM}},
\kappa_{\mathrm{fiber}})=(0,0,3,5,24)\) at \(K3\times E\); the row
\(\kappa_{\mathrm{BKM}}=5\) is the cusp-form weight realized by
\(\Delta_5\), and the row \(\kappa_{\mathrm{fiber}}=24\) is
\(\chi_{\mathrm{top}}(K3)=24\).  The additive expression
\(\kappa_{\mathrm{BKM}}=\kappa_{\mathrm{ch}}+\chi(\mathcal O_{\mathrm{fiber}})\)
has no invariant meaning: the Hodge reading gives \(0+2=2\) at \(N=1\),
and the Heisenberg reading gives the accidental equality \(3+2=5\) only
on that row.  The CHL frame has Borcherds weights
\((5,3,2,\tfrac32,1)\), fixed by \(c_N(0)/2\) with
\(c_N(0)=(10,6,4,3,2)\) per Jatkar--Sen and Govindarajan--Krishna
\cite{JS06Dyon,GK09CHL,GK10Composite},
not by a fixed additive
formula.  Thus \(\kappa_{\mathrm{fiber}}\) is the Euler characteristic
of the topological fibre and not of \(\mathcal O_{\mathrm{fiber}}\); for
K3, \(\chi_{\mathrm{top}}(K3)=24\) but
\(\chi(\mathcal O_{K3})=2\).  Under
the Vol III chain-fusion datum at level \(2\), the chiral algebra
\(A_X\) on the CY face is compared with the chart algebra
\(A_b\) on the Vol II open face.  In the \(K3\times E\) specialization
this comparison becomes
\(A_X=A_b=\Omega_E^{\mathrm{ch}}C_X\simeq\mathsf{Den}_{\Delta_5,E}\) only
under the hypotheses of Theorem~\ref{thm:ch6-primitive-recognition}.

\begin{remark}[Coefficient projection is not a Hall comparison]
\label{rmk:ch6-coefficient-projection}
A Hall--Borcherds comparison map cannot be defined by coefficient
projection alone, and a coefficient projection respecting charge
addition because both products use the same cone is not a recognition
theorem.  A coefficient projection is not a Hall product/bracket matrix
comparison, and does not prove Borcherds--Serre rows, radical descent,
no-extra-relations, or PBW.  The recognition criterion (S4)--(S10) of
Theorem~\ref{thm:ch6-primitive-recognition} is the mathematical
statement:
the source-side bracket matrices, Hopf radical, kernel equality, and
PBW associated-graded comparison are independent data, none of which is
inferable from a coefficient projection.
\end{remark}

\subsection{BCOV / mixed holomorphic-topological strings}
\label{subsec:ch6-bcov}

The mixed holomorphic-topological strings companion provides the local
model \(\R^2_{\mathrm{top}}\times\C^2_{\mathrm{hol}}\) with brane
completion as one realization of the factorization-algebra datum
\(\Phi_d^{\mathrm{FA}}\); the global obstruction is the holomorphic
de Rham obstruction class.  The level-\(n\) scalar shadow
\(Z_{\mathrm{BPS}}^{K3\times E}=\Delta_5^{-2}\) on the closed
\(K3\times E\mapsto\mathrm{point}\) cobordism is the BPS partition
function (Chapter~\ref{ch:zbps-realization}).  A three-dimensional
gravity-line interpretation with internal compact factor \(K3\times E\)
requires the Hall--Borcherds residual of Vol II and is not claimed by
Theorem~\ref{thm:ch6-pfaffian-dirac}.  The identity
\(\operatorname{Pf}_{\mathrm{prot}}(\mathfrak D_X^{\mathrm{DI}})=\mathcal D_X=\Delta_5\)
is a protected Pfaffian identity before trace; its inverse square is
the level-\(n\) protected scalar trace.  The Pfaffian and orientation
parts of the operator-level datum \(\mathfrak D_X^{\mathrm{DI}}\) are
defined on \(K3\times E\)
relative to the retained data of
Appendix~\ref{app:obstruction-discharges}; the primitive Lie-superalgebra
part is the compact-source recognition problem.

\section{Residual obstructions and open subroutes}
\label{sec:ch6-residuals}

The Pfaffian and orientation theorems are retained-data theorems on
\(K3\times E\), relative to Appendix~\ref{app:obstruction-discharges}.
The residuals in this
chapter are the compact-source recognition datum \((S1)\)--\((S10)\)
and the gravity-line Hall--Borcherds operator algebra.

\paragraph{(D0).}  Theorem~\ref{thm:G-D0} is the criterion saying that
the retained D0-HN datum supplies the cosection-twisted reduced
d-critical orientation through the D0-degeneration pro-limit and the
cofinal HN inverse system.  The \(D_0\)-state, \(D_0\)-quotient,
\(D_0\)-observable, \(D_0\)-descent, and \(D_0\)-Pfaffian rows are
therefore part of that datum.

\paragraph{(O1).}  Theorem~\ref{thm:G-O1} supplies the strong reduced
orientation on the \(K3\times E\)-quotient self-Ext frame bundle
relative to the retained quotient-orientation datum.  The
finite-stabilizer linearization characters and the connected
\(BE\)-edge class are part of that datum, not consequences of the
primitive recognition theorem.

\paragraph{(O1)\(^+\).}  The equivariance part of
Theorem~\ref{thm:G-O1-plus} is relative to
the chosen Weyl lifts, the vanishing projective cocycle, and the zero
torsor defects on the retained groupoid.  It gives the honest
\(W^{(2)}(\La^{2,1}_{II})\)-transport of the
orientation character.  The comparison with \(\nu_{\Delta_5}\) also
uses the rank-one wall signs from \((O2_{\mathrm{atlas}})\).  Neither
statement supplies the relation-closed source matrices of the
recognition datum.

\paragraph{(O2).}  Relative to the retained type-II wall-atlas datum
\((O2_{\mathrm{atlas}})\), the rank-one normal-form rank
\(N_{\delta_i}^{\mathrm{Pf}}=1\) holds for \(i=1,2,3\)
(Proposition~\ref{prop:appE-local-pfaffian-sign}).  The type-II middle-wall
constant \(4096=2^{12}=64^2\) admits two readings: as the squared
theta-leading on the type-II Borcherds product (where the \(64=2^6\)
prefactor of \(D_5=64^{-1}\Delta_5\) doubles by \(\Delta_5^2=\Delta_{10}\))
and as the \(2^{12}\) sign assignments over the twelve binary wall
coordinates of the half-Hilbert orbit.  The equality of these two
readings is scalar only, by Theorem~\ref{thm:G-R1-R2}.  The
OP scalar branch sign
\(-1\) is recorded by \((-p_{\mathrm{DT}})^n\) of the OP/DT chamber,
not by the orientation line.  The middle wall
\(\delta_2\leftrightarrow(0,1,1)\) is represented by a wrapped wall
object with positive geometric \(b\)-degree only after the retained
\((O2_{\mathrm{atlas}})\) datum supplies that object.
Proposition~\ref{prop:appE-local-pfaffian-sign}
is the rank-one local sign calculation; the component, boundary,
overlap, quotient, and protected-integration compatibilities are the
global content of \((O2_{\mathrm{atlas}})\).  Higher-rank wall atlases
are not used in the proof of Theorem~\ref{thm:ch6-pfaffian-dirac}.

\paragraph{Recognition criterion (S1)--(S10).}  The window
\(W_{\le 3}\) is the first arithmetic window but is not relation-closed.
Put \(a_{ij}=\delta_i+\delta_j\),
\(\delta_{123}=\delta_1+\delta_2+\delta_3\), and
\(C_{k,s}=a_{ij}+s\delta_k\) for \(\{i,j,k\}=\{1,2,3\}\).
The first base-string terminal degree window is the height-\(7\) closure
\[
W_{\mathrm{rel}}^+
=W_{\le 3}^{\mathrm{act}}\cup R_4\cup I_4\cup C_{\le 7}
\cup\{2\delta_{123}\},
\]
where \(R_4=\{3\delta_i+\delta_j:i\ne j\}\),
\(I_4=\{2a_{ij}:i<j\}\), and
\(C_{\le7}=\{C_{k,s}:2\le s\le5\}\).
It adds the real-Serre terminal rows, doubled isotropic rows,
\(\delta_{123}\)-real strings, odd--odd \(\delta_{123}\) rows, and
complementary height-seven strings.  In this enlarged window,
the finite target degree-closure packet
\(\texttt{wrel\_degree\_closure}\) checks the \(35\) rows, their
\(\delta\)-basis coordinates, heights, norms, Jacobi signed
coefficients, six real-Serre terminal zero rows, and the three
complementary-string terminal zero rows.  Its downward-saturation
audit proves the precise defect: \(W_{\mathrm{rel}}^+\) is not
downward saturated, and the only nonzero proper subdegrees outside
the \(35\)-row set are
\[
D_k=\delta_k+2a_{ij},\qquad \{i,j,k\}=\{1,2,3\},
\]
all below \(2\delta_{123}\), with signed multiplicity
\(\smult(D_k)=-513\).
The same packet verifies that adjoining these three rows gives the
minimal nonzero downward-saturation extension
\[
W_{\mathrm{sat}}^+
=W_{\mathrm{rel}}^+\cup\{D_1,D_2,D_3\},
\]
with zero remaining nonzero proper-subdegree defect.
Remark~\ref{rem:bkm-relation-window-target-parity} supplies full target
parity for \(I_4=\{2a_{ij}\}\), \(C_{k,3}\), \(C_{k,4}\), and
\(C_{k,5}\), while \(C_{k,2}\), \(D_k\), and \(2\delta_{123}\) remain
signed target rows until the finite target presentation is computed.
The finite coefficient table
\(\texttt{k3e\_first\_relation\_window\_coefficients}\) checks the
Jacobi arithmetic for these rows:
\[
f(4,4)=10,\qquad f(2,2)=108,\qquad
f(2,1)=f(4,3)=-513,\qquad
f(3,3)=-64,\qquad f(4,2)=4016,
\]
with the terminal rows \(f(5,5)=f(1,-3)=0\).  This is a target
signed-dimension check, not a parity split or a source comparison.
The A071 target-presentation verifier separately checks the target
parity rows \(2a_{ij}:10|0\), \(C_{k,3}:29|93\),
\(C_{k,4}:10|0\), and \(C_{k,5}:0|0\), and records
\(C_{k,2}\), \(D_k\), and \(2\delta_{123}\) as signed-only rows with
no target basis, pairing, radical, or PBW blocks.
Proposition~\ref{prop:target-admissibility-obstruction-wrel} proves
that these signed rows have no admissible parity, pairing, radical, or
PBW target blocks for (R0)--(R8).  Thus \(W_{\mathrm{rel}}^+\) is a
base-string terminal degree set with a three-dimensional saturation
defect, and \(W_{\mathrm{sat}}^+\) is a target-recorded saturated
extension with signed-only rows.  The post-saturation real-string audit
then finds \(21\) terminal codomains outside \(W_{\mathrm{sat}}^+\),
ten with nonzero signed multiplicity.  Full finite relation closure
therefore starts beyond \(W_{\mathrm{sat}}^+\).  The terminal-layer
audit groups these codomains into \(21\) distinct terminal degrees;
adjoining them creates \(114\) nonzero subdegree-defect incidences
across \(26\) distinct subdegrees.  The next target-degree computation
adjoins exactly these \(26\) rows and gives an \(85\)-row degree set with
zero remaining nonzero proper-subdegree defect.  This closes downward
saturation for the displayed terminal layer, not the real-string
relation audit and not the compact-source comparison.  If these \(26\)
rows are promoted to target-generator rows, the real-root strings force
\(55\) terminal checks; \(54\) terminal codomains lie outside the
\(85\)-row set, and \(12\) of the missing codomains have nonzero signed
multiplicity.  If those \(55\) terminal codomains are adjoined, the next
downward-saturation audit already has \(546\) nonzero proper-subdegree
defect incidences across \(98\) distinct subdegrees.  Adjoining those
\(98\) rows gives a \(237\)-row target degree set with zero remaining
nonzero proper-subdegree defect for that layer.  If the \(98\)-row layer
is promoted to target-generator rows, the real-root strings force
\(206\) terminal checks, \(182\) new terminal degrees, and \(24\)
missing nonzero terminal degrees.  The exact terminal-layer table
records those \(206\) codomains, and the exact saturation-extension
table records \(50\) nonzero proper subdegrees with \(624\) incidences;
adjoining both gives a \(469\)-row target degree set with zero remaining
nonzero proper-subdegree defects for that layer.  Promoting the
\(50\)-row layer forces \(96\) real-string checks, \(90\) new terminal
degrees, no missing nonzero terminal degree, and a \(12\)-row saturation
extension with \(44\) incidences; adjoining both gives a \(571\)-row
target degree set with zero remaining nonzero proper-subdegree defects
for that layer.  Promoting the \(12\)-row layer forces \(20\) zero
terminal checks and a \(6\)-row saturation extension with \(8\)
incidences; adjoining both gives a \(597\)-row target degree set with
zero remaining nonzero proper-subdegree defects for that layer.
Promoting the newly adjoined \(6\)-row layer forces \(10\) zero
terminal checks and a \(6\)-row saturation extension with \(18\)
incidences; adjoining both gives a \(613\)-row target degree set with
zero remaining nonzero proper-subdegree defects for that layer.
Promoting the next newly adjoined \(6\)-row layer forces \(10\) zero
terminal checks and a \(6\)-row saturation extension with \(8\)
incidences; adjoining both gives a \(629\)-row target degree set with
zero remaining nonzero proper-subdegree defects for that layer.
Promoting the next newly adjoined \(6\)-row layer forces \(10\) zero
terminal checks and a \(6\)-row saturation extension with \(18\)
incidences; adjoining both gives a \(645\)-row target degree set with
zero remaining nonzero proper-subdegree defects for that layer.  The
\(\texttt{wrel\_tail\_staircase}\) verifier then proves the symbolic
tail \(A_N\Rightarrow B_N\Rightarrow A_{N+2}\) for every even
\(N\ge14\), and its finite-anchor table checks \(A_{14}\) inside the
\(\texttt{wrel\_degree\_closure}\) packet.  Hence the target-degree
audit has an infinite six-row staircase; its unbounded-tail table
records a strictly increasing coordinate subsequence.  No finite prefix
of this target-only audit is a primitive-recognition theorem.  Neither
displayed target window is yet a primitive-recognition
window satisfying (R0)--(R8).  At
\(W_{\le 3}\)
the target parity
table
\(\delta_i:1|0\), \(a_{ij}:10|0\), \(2\delta_i+\delta_j:1|0\),
\(\delta_{123}:29|93\) is computed; the row \(a_{ij}:10|0\) is the
real bracket \([e_i,e_j]\) together with the nine Gritsenko--Nikulin
isotropic simple directions.  The
\(\texttt{wle3\_target\_parity}\) verifier checks this target table and
its scalar firewall; a compact source implementing
(S3)--(S10) on \(W_{\le 3}\) is open, and the source identification at
\(W_{\mathrm{rel}}^+\) requires a finite source-and-target
presentation computation.
The first relation-closed source theorem is exactly the finite packet
\[
\mathrm{window\_closure},\quad
\mathrm{target\_source\_representatives},\quad
\mathrm{chevalley\_serre\_matrices},\quad
\mathrm{hall\_boundary\_complex},\quad
\mathrm{spectral\_sequence},\quad
\mathrm{kernel\_matrices},\quad
\mathrm{pbw\_graded},\quad
\mathrm{first\_window\_theorems},\quad
\mathrm{transitions},\quad
\mathrm{scalar\_firewall}.
\]
It must compute the compact representatives, Chevalley and Serre
matrices, Hall--Chevalley boundary, kernel bases, no-extra equality,
PBW associated gradeds, and first-window theorem rows.  It is not
provided by the target parity table, the signed Borcherds exponents,
the Pfaffian product, or the scalar trace.
The current packet
\(\texttt{certificates/first\_window/k3e\_relation\_closed\_window}\)
records these missing rows by the obstruction ledger
\(\texttt{FIRST\_WINDOW\_OBSTRUCTION\_LEDGER\_VERIFIED}\), with
\(\texttt{first\_window\_certification=false}\).  Hence Chapter~\ref{ch:pfaffian-dirac}
uses no first-window primitive-recognition theorem beyond this
checked absence status.

\paragraph{Vol III Hall--Drinfeld--Borcherds spectral sequence on compact
non-toric CY3.}  The Hall--Drinfeld--Borcherds spectral sequence
abutting from Vol III's CoHA-type algebra to the BKM denominator on
\(K3\times E\) is open on the compact non-toric trivial-canonical threefold,
even after (D0), (O1), (O1)\(^+\), and (O2).  The
finite-stage Hall source kernel of
Theorem~\ref{thm:finite-stage-dirac-igusa-pro-object} is the
finite-stage substitute.

\paragraph{Protected determinant and character level.}
The Pfaffian and orientation theorems identify the
\textcircled{\scriptsize 2}\(\leftrightarrow\)\textcircled{\scriptsize 3}
edge at the protected determinant and character level on \(K3\times E\).
The primitive Lie-superalgebra recognition remains the finite
compact-source datum \((S1)\)--\((S10)\).  The
cross-level identities
\textcircled{\scriptsize 1}\(\leftrightarrow\)\textcircled{\scriptsize 2}
and \textcircled{\scriptsize 2}\(\leftrightarrow\)\textcircled{\scriptsize 5}
are imported from Borcherds--Gritsenko--Nikulin and the regularized
theta lift.  The
operator-level datum
\(\mathfrak D_X^{\mathrm{DI}}\) is defined at the Pfaffian and
orientation level relative to the retained data of
Appendix~\ref{app:obstruction-discharges}; the
scalar shadow \(Z_{\mathrm{BPS}}^{K3\times E}=\Delta_5^{-2}\) is the
level-\(n\) protected trace.  The compact microscopic state-space
interpretation requires the compact-source recognition datum; the
\(3d\) gravitational path integral requires the Hall--Borcherds
residual; the Calabi--Yau realization requires a
construction external to the Pfaffian theorem.
